\section{Introduction}
\label{sec:introduction}

Users access a variety of free content and services on the web, which are largely funded through online advertising.
%
In turn, advertising is heavily dependent on monitoring users' online activities for various purposes such as analytics, (re)targeting, and conversion tracking.
%
To realize these objectives, user tracking has become a pervasive part of the web---today, 92\% of webpages embed at least one third-party tracker~\cite{urban2024WebAlmanac2024}.

%%
Over the years, researchers have conducted numerous studies to examine the evolution of web tracking mechanisms, browser developments, and regulatory compliance.
%
However, despite a considerable body of work, major findings remain scattered across many disparate studies, and existing survey papers have either been too high-level~\cite{ermakova2018web}, too narrowly-focused on specific forms of web tracking
such as caching and browser fingerprinting
~\cite{bujlowSurveyWebTracking2017,laperdrix2020browser}, or are not current with modern techniques anymore~\cite{mayerThirdPartyWebTracking2012}.
%
Furthermore, as privacy defenses improve in browsers, trackers continually adapt with new evasion techniques~\cite{narayanan2018web,localmess-usenix-sec-26}.
%
The result is an ever-shifting technical landscape of tracking techniques.

%%
Regulations often govern tracking practices and although regulatory changes have had a more gradual impact than browser-based technical interventions, together they have continued to reshape the ecosystem.
%
Today, web tracking is undergoing a transformative change due to the introduction of privacy-enhancing protections in major web browsers, new advertising paradigms that do not rely on cross-site tracking, and evolving regulatory frameworks~\cite{hercherW3CAdPrivacy2022,chavezUpdatePlansPrivacy2025}.
%
These shifts call for a comprehensive study of emerging trends in the evolving tracking landscape to identify crucial research gaps.

%%
To this end, we synthesize disparate lines of research and practices in web tracking---spanning across technical mechanisms, mitigations, and regulatory changes---to systematically provide an overview of the current state of web tracking.
%
We scope this SoK to how the data is \textit{collected} about users, not how that data might then be \textit{used}.
%
This unified perspective enables a critical reflection on how far the community has come and which research directions it should prioritize.
%
Our main contributions are as following:
%
\begin{itemize}
    \item We systematize the extensive body of research on web tracking, providing a consolidated knowledge base of advances, highlighting evolving trends, and bridging emerging but related tracking mechanisms.
    \item We systematically organize anti-tracking interventions proposed by researchers, major browser providers, and community as well as relevant regulatory and compliance actions across the EU and the US to assess how they have altered the ecosystem over the years.
    \item We identify key open challenges and promising future research directions in web tracking.
\end{itemize}
